\documentclass[10pt]{article}

\usepackage[utf8]{inputenc}
\usepackage[T1]{fontenc}
\usepackage{amsmath,amssymb,mathtools}
\usepackage{booktabs}
\usepackage{enumitem}
\usepackage[margin=0.78in]{geometry}
\usepackage{xcolor}
\usepackage[most]{tcolorbox}
\usepackage{listings}
\usepackage{hyperref}

\definecolor{ResultBlue}{RGB}{31,78,121}
\definecolor{ProjectGold}{RGB}{190,128,20}
\hypersetup{colorlinks=true,linkcolor=ResultBlue,urlcolor=ResultBlue}
\setlist{itemsep=2pt,topsep=3pt,parsep=1pt,partopsep=0pt}
\setlength{\parindent}{0pt}
\setlength{\parskip}{4pt}
\setlength{\abovedisplayskip}{5pt plus 2pt minus 2pt}
\setlength{\belowdisplayskip}{5pt plus 2pt minus 2pt}

\newcommand{\R}{\mathbb{R}}
\newcommand{\trans}{^{\mathsf T}}
\newcommand{\norm}[1]{\left\lVert#1\right\rVert}
\newcommand{\diag}{\operatorname{diag}}
\newcommand{\proj}{\operatorname{proj}}
\newcommand{\pospart}[1]{\left(#1\right)_{+}}

\newtcolorbox{resultbox}[1][Result]{breakable,colback=ResultBlue!3,
  colframe=ResultBlue,boxrule=0.8pt,arc=1mm,
  title=\textbf{#1},fonttitle=\bfseries}
\newtcolorbox{projectbox}[1]{breakable,colback=ProjectGold!7,
  colframe=ProjectGold!90!black,boxrule=0.8pt,arc=1mm,
  title=\textbf{#1},fonttitle=\bfseries}

\lstdefinestyle{python}{
  language=Python,
  basicstyle=\ttfamily\scriptsize,
  keywordstyle=\color{ResultBlue}\bfseries,
  commentstyle=\color{black!55},
  stringstyle=\color{ProjectGold!80!black},
  showstringspaces=false,
  breaklines=true,
  columns=fullflexible,
  frame=single,
  rulecolor=\color{black!25},
  xleftmargin=2pt,
  xrightmargin=2pt
}

\title{Summer Bootcamp -- Session 3\\Corrections to the Optimization Exercises}
\author{Nathan Sauldubois}
\date{August 2026}

\begin{document}
\maketitle

\section*{Correction Guidelines}

The solutions follow the numbering of the exercise sheet. Necessary conditions
are not used as sufficient conditions unless convexity, curvature, or a direct
comparison supplies the missing argument. Numerical experiments use the seeds
stated below and stop when
\(\norm{\nabla f(x_k)}_2\leq10^{-8}\), or, for a constrained problem, when
both feasibility and stationarity residuals are at most \(10^{-8}\).

The companion file
\texttt{session\_3\_optimization\_corrections.py} contains the complete
reference computations for Exercises 10 and 11. The separate Portfolio
Optimization Project is not corrected in this document; its status is stated
at the end.

% -----------------------------------------------------------------------------
\section{One-Dimensional Unconstrained Optimization}

\paragraph{Exercise 1 -- Critical points are only candidates.}

For \(f_1(x)=x^4-4x^2+3\),
\[
f_1'(x)=4x(x^2-2),\qquad f_1''(x)=12x^2-8.
\]
Thus \(0\) is a strict local maximum and \(\pm\sqrt2\) are strict local
minima. Since \(f_1(x)\to+\infty\) as \(|x|\to\infty\), the two minima are
global:
\[
f_1(0)=3,\qquad f_1(\pm\sqrt2)=-1.
\]

For \(f_2(x)=x^3-3x+1\), the interior critical points are \(\pm1\), with
\(f_2''(x)=6x\). The boundary comparison is essential:
\[
\begin{array}{c|rrrr}
x&-2&-1&1&2\\ \midrule
f_2(x)&-1&3&-1&3
\end{array}
\]
Hence the global minima on \([-2,2]\) are \(x=-2\) and \(x=1\), while the
global maxima are \(x=-1\) and \(x=2\). Among the interior points, \(-1\) is a
strict local maximum and \(1\) a strict local minimum. Relative to the closed
domain, the two endpoints are also one-sided local extrema.

Finally, for
\(f_3(x)=\log(1+e^x)-\frac34x\),
\[
f_3'(x)=\frac{e^x}{1+e^x}-\frac34,
\qquad
f_3''(x)=\frac{e^x}{(1+e^x)^2}>0.
\]
The function is strictly convex, and stationarity gives \(e^x=3\). Moreover,
\(f_3(x)\to+\infty\) at both ends of the real line. Therefore
\begin{resultbox}
\[
x_3^\star=\log3,
\qquad
f_3(x_3^\star)=\log4-\frac34\log3,
\]
and this minimizer is unique and global.
\end{resultbox}
The sketches must show the symmetric double well of \(f_1\), the cubic on its
restricted interval for \(f_2\), and the strictly convex bowl of \(f_3\), with
the points just identified marked explicitly.

\paragraph{Exercise 2 -- A parameter changes the number of local optima.}

Stationarity for
\(f_a(x)=\frac14x^4-\frac a2x^2+x\) is
\[
p_a(x):=f_a'(x)=x^3-ax+1=0.
\]
The turning points of this cubic occur at \(x=\pm s\), where
\(s=\sqrt{a/3}\), and
\[
p_a(-s)=1+\frac{2a^{3/2}}{3\sqrt3}>0,
\qquad
p_a(s)=1-\frac{2a^{3/2}}{3\sqrt3}.
\]
Introduce the threshold
\[
a_c=\left(\frac{3\sqrt3}{2}\right)^{2/3}
=\frac{3}{2^{2/3}}.
\]
Then \(p_a\) has one real root for \(0<a<a_c\), three distinct real roots for
\(a>a_c\), and a simple root plus a double root when \(a=a_c\).

In the three-root regime, write the roots \(r_1<r_2<r_3\). The sign pattern is
\[
\begin{array}{c|ccccccc}
x&(-\infty,r_1)&r_1&(r_1,r_2)&r_2&(r_2,r_3)&r_3&(r_3,\infty)\\ \midrule
f_a'(x)&-&0&+&0&-&0&+
\end{array}
\]
so \(r_1,r_3\) are strict local minima and \(r_2\) is a strict local maximum.
At \(a=a_c\), the negative simple root is a strict local minimum; the positive
double root of \(p_a\) is stationary but is neither a minimum nor a maximum.

Because \(f_a(x)\to+\infty\) as \(|x|\to\infty\), a global minimum exists.
Checking \(f_a'=0\) only generates its candidates. A rigorous comparison can
evaluate \(f_a\) at all stationary points; at a stationary root \(r\), the
identity \(r^3=ar-1\) simplifies the value to
\[
f_a(r)=-\frac a4r^2+\frac34r.
\]
There is also a direct argument in the three-root regime: for every \(x>0\),
\(f_a(-x)=f_a(x)-2x<f_a(x)\). Hence the positive local minimum cannot be
global, and the negative local minimum \(r_1\) is the unique global minimizer.

\paragraph{Exercise 3 -- Newton's method as a local model.}

Here
\[
f'(x)=e^x-3+\frac x2,
\qquad
f''(x)=e^x+\frac12>0.
\]
Strict convexity and coercivity give one global minimizer. It is the unique
root of \(f'\), and it may be written using the Lambert \(W\) function as
\[
x^\star=6-W(2e^6).
\]
Newton's iteration is
\begin{resultbox}
\[
x_{k+1}=x_k-
\frac{e^{x_k}-3+x_k/2}{e^{x_k}+1/2}.
\]
\end{resultbox}
Indeed, the quadratic model in a displacement \(p\) is
\[
m_k(p)=f(x_k)+f'(x_k)p+\frac12f''(x_k)p^2,
\]
whose unique minimizer is \(p_k=-f'(x_k)/f''(x_k)\). At every nonstationary
\(x_k\),
\[
f'(x_k)p_k=-\frac{f'(x_k)^2}{f''(x_k)}<0,
\]
so \(p_k\) is a descent direction. This does not guarantee that the full step
\(\alpha_k=1\) decreases the true nonlinear objective. Armijo backtracking
replaces it by \(x_{k+1}=x_k+\alpha_kp_k\), repeatedly reducing \(\alpha_k\)
until
\[
f(x_k+\alpha_kp_k)
\leq f(x_k)+c\alpha_k f'(x_k)p_k.
\]
It therefore preserves the fast local Newton step when it is acceptable and
forces sufficient decrease when the quadratic model is poor.

% -----------------------------------------------------------------------------
\section{Multidimensional Unconstrained Optimization}

\paragraph{Exercise 4 -- Quadratic optimization and conditioning.}

For \(f(x)=\frac12x\trans Qx-c\trans x\), with \(Q=Q\trans\succ0\),
\[
\nabla f(x)=Qx-c,\qquad \nabla^2f(x)=Q.
\]
For distinct \(x,y\) and \(t\in(0,1)\), direct expansion gives
\[
tf(x)+(1-t)f(y)-f(tx+(1-t)y)
=\frac12t(1-t)(x-y)\trans Q(x-y)>0,
\]
so \(f\) is strictly convex. Positive definiteness also implies
\(\ker Q=\{0\}\), hence \(Q\) is invertible. The stationarity system has the
unique solution
\[
\boxed{x^\star=Q^{-1}c},
\]
where a numerical implementation should use a linear solve. Completing the
square yields the global certificate
\[
f(x)-f(x^\star)=\frac12(x-x^\star)\trans Q(x-x^\star)\geq0.
\]

Let \(Q=U\diag(\lambda_1,\ldots,\lambda_d)U\trans\). In coordinates
\(z=U\trans(x-x^\star)\), the objective gap is
\(\frac12\sum_i\lambda_i z_i^2\): the eigenvectors are principal directions
and the eigenvalues are their curvatures.

For gradient descent, \(e_k=x_k-x^\star\) satisfies
\[
e_{k+1}=(I-\alpha Q)e_k,
\qquad
(U\trans e_{k+1})_i=(1-\alpha\lambda_i)(U\trans e_k)_i.
\]
Convergence for every initial point is equivalent to
\[
0<\alpha<\frac{2}{\lambda_{\max}}.
\]
When \(\kappa(Q)\approx1\), level sets are nearly spherical and all
coordinates contract at comparable rates. When \(\kappa(Q)\gg1\), the level
sets are elongated: a step small enough for the largest curvature makes slow
progress along the smallest-curvature direction. The constant step
\(2/(\lambda_{\min}+\lambda_{\max})\) has worst-coordinate contraction factor
\((\kappa-1)/(\kappa+1)\).

\paragraph{Exercise 5 -- Least squares, ridge regression, and conditioning.}

For
\(\ell(x)=\frac12\norm{Ax-b}_2^2\),
\[
\nabla\ell(x)=A\trans(Ax-b),
\]
so every minimizer solves the normal equations
\[
A\trans A x=A\trans b.
\]
The solution is unique exactly when \(\ker A=\{0\}\), equivalently when \(A\)
has full column rank. At any solution \(x^\star\),
\[
A\trans(b-Ax^\star)=0,
\]
so the residual is orthogonal to the column space of \(A\); \(Ax^\star\) is
the orthogonal projection of \(b\) onto that space.

For ridge regression,
\[
\nabla\ell_\lambda(x)
=(A\trans A+\lambda I)x-A\trans b,
\]
and therefore
\begin{resultbox}
\[
x_\lambda=(A\trans A+\lambda I)^{-1}A\trans b.
\]
\end{resultbox}
For every nonzero \(v\),
\[
v\trans(A\trans A+\lambda I)v
=\norm{Av}_2^2+\lambda\norm v_2^2>0.
\]
Thus ridge has a unique solution even if \(A\) is rank deficient.

If \(\sigma_i\) are the singular values of \(A\), the Hessian eigenvalues
become \(\sigma_i^2+\lambda\). Hence
\[
\kappa(A\trans A+\lambda I)
=\frac{\sigma_{\max}^2+\lambda}{\sigma_{\min}^2+\lambda},
\]
where \(\sigma_{\min}=0\) in a rank-deficient problem. Regularization raises
small curvatures and improves stability, but the singular-vector coefficients
are shrunk by factors \(\sigma_i/(\sigma_i^2+\lambda)\). This creates bias in
exchange for lower sensitivity to noise.

\paragraph{Exercise 6 -- Second-order conditions in several dimensions.}

For \(f_1\),
\[
\nabla f_1=\binom{2x+y-4}{x+2y+2},
\qquad
H_1=\begin{pmatrix}2&1\\1&2\end{pmatrix}.
\]
The stationary point is
\((10/3,-8/3)\), and the Hessian eigenvalues \(1,3\) are positive. It is the
unique global minimizer.

For \(f_2\), stationarity requires
\[
4x(x^2-1)=0,
\qquad
4y(y^2-2)=0,
\]
so \(x\in\{0,\pm1\}\) and \(y\in\{0,\pm\sqrt2\}\). Moreover,
\[
H_2(x,y)=\diag(12x^2-4,12y^2-8).
\]
The classification is
\[
\begin{array}{c|c|c}
\text{points}&\text{Hessian signs}&\text{classification}\\ \midrule
(\pm1,\pm\sqrt2)&(+,+)&\text{strict global minima, value }-5\\
(0,0)&(-,-)&\text{strict local maximum, value }0\\
(\pm1,0)&(+,-)&\text{saddles}\\
(0,\pm\sqrt2)&(-,+)&\text{saddles}
\end{array}
\]
There are no degenerate stationary candidates for the stated \(f_2\): both
diagonal Hessian entries are nonzero at every candidate. Thus the requested
directional inspection is not needed here; this observation is itself the
appropriate check.

For \(f_3\), the only stationary point is \((0,0)\), and
\[
H_3=\begin{pmatrix}2&2\\2&-2\end{pmatrix},
\qquad \det H_3=-8<0.
\]
It is a saddle. The restrictions \(f_3(x,0)=x^2\) and
\(f_3(0,y)=-y^2\) give the same conclusion and show that the function is
unbounded above and below. Contour plots should respectively display an
elliptic bowl, four equivalent wells for \(f_2\), and hyperbolic saddle
contours for \(f_3\).

% -----------------------------------------------------------------------------
\section{Multidimensional Constrained Optimization}

\paragraph{Exercise 7 -- Projection onto an affine constraint.}

The projection is the solution of
\[
\min_x\ \frac12\norm{x-z}_2^2
\quad\text{subject to}\quad a\trans x=b.
\]
With
\(\mathcal L(x,\lambda)=\frac12\norm{x-z}^2+\lambda(a\trans x-b)\),
stationarity gives \(x=z-\lambda a\). Enforcing feasibility yields
\begin{resultbox}
\[
\proj_C(z)=z-\frac{a\trans z-b}{\norm a_2^2}a.
\]
\end{resultbox}
The correction vanishes when \(z\in C\), and
\(z-\proj_C(z)\in\operatorname{span}(a)\), the normal space of the hyperplane.

For \(C=\{x:Ax=b\}\), use the vector multiplier \(\lambda\). Stationarity and
feasibility give
\[
x=z-A\trans\lambda,
\qquad
AA\trans\lambda=Az-b.
\]
If \(A\) has full row rank, then \(A\trans y=0\) implies \(y=0\), and
\(y\trans AA\trans y=\norm{A\trans y}^2>0\) for \(y\ne0\). Thus \(AA\trans\)
is invertible and
\[
\boxed{\proj_C(z)=z-A\trans(AA\trans)^{-1}(Az-b).}
\]

\paragraph{Exercise 8 -- Equality constraints, duality, and a saddle point.}

Set \(S=AQ^{-1}A\trans\). The Lagrangian and KKT system are
\[
\mathcal L(x,\lambda)
=\frac12x\trans Qx+c\trans x+\lambda\trans(Ax-b),
\]
\[
\begin{pmatrix}Q&A\trans\\A&0\end{pmatrix}
\binom{x^\star}{\lambda^\star}
=\binom{-c}{b}.
\]
Since \(Q\succ0\) and \(A\) has full row rank, \(S\succ0\). Eliminating \(x\)
gives
\begin{resultbox}
\[
\lambda^\star=-S^{-1}(b+AQ^{-1}c),
\]
\[
x^\star=-Q^{-1}c
+Q^{-1}A\trans S^{-1}(b+AQ^{-1}c).
\]
\end{resultbox}

For a fixed multiplier, minimizing the Lagrangian in \(x\) at
\(x(\lambda)=-Q^{-1}(c+A\trans\lambda)\) gives
\[
q(\lambda)
=-\frac12(c+A\trans\lambda)\trans Q^{-1}
  (c+A\trans\lambda)-b\trans\lambda.
\]
Its Hessian is \(-S\prec0\), so the dual function is strictly concave.

If \(x\) is feasible, then for every \(\lambda\),
\[
q(\lambda)=\inf_y\mathcal L(y,\lambda)
\leq\mathcal L(x,\lambda)=f(x),
\]
which is weak duality. At the KKT pair, \(x^\star\) minimizes
\(\mathcal L(\cdot,\lambda^\star)\) and is feasible, hence
\[
q(\lambda^\star)=\mathcal L(x^\star,\lambda^\star)=f(x^\star).
\]
This proves strong duality and the saddle relations
\[
\mathcal L(x^\star,\lambda)
=\mathcal L(x^\star,\lambda^\star)
\leq\mathcal L(x,\lambda^\star)
\quad\text{for all }x,\lambda.
\]
Finally, if \(V(b)\) is the optimal value, the envelope theorem with the
stated sign convention gives
\[
DV(b)[\Delta b]=-\lambda^{\star\trans}\Delta b,
\qquad
\nabla_bV(b)=-\lambda^\star.
\]

\paragraph{Exercise 9 -- Projection onto a half-space with KKT.}

For the inequality \(g(x)=a\trans x-b\leq0\),
\[
\mathcal L(x,\lambda)
=\frac12\norm{x-z}^2+\lambda(a\trans x-b).
\]
The KKT conditions are
\[
x-z+\lambda a=0,
\qquad a\trans x\leq b,
\qquad \lambda\geq0,
\qquad \lambda(a\trans x-b)=0.
\]
If \(a\trans z\leq b\), the unconstrained minimizer is feasible and
\((x^\star,\lambda^\star)=(z,0)\). If \(a\trans z>b\), the constraint is active
and
\[
\lambda^\star=\frac{a\trans z-b}{\norm a^2},
\qquad
x^\star=z-\frac{a\trans z-b}{\norm a^2}a.
\]
Consequently,
\begin{resultbox}
\[
\proj_{\{a\trans x\leq b\}}(z)
=z-\frac{\pospart{a\trans z-b}}{\norm a_2^2}a.
\]
\end{resultbox}
The positive part switches continuously from an inactive to an active
constraint at the boundary. The objective is strictly convex, the constraint
is affine, and a strictly feasible point exists because \(a\ne0\); therefore
the KKT conditions are sufficient (and the optimizer is unique).

% -----------------------------------------------------------------------------
\section{Python Algorithms}

\paragraph{Exercise 10 -- Implement gradient descent from scratch.}

A reusable implementation must evaluate the stopping test before taking the
next step and must keep the initial point in its history. The essential loop is
shown below; the companion Python file adds input checks and all requested
plots.

\begin{lstlisting}[style=python]
def gradient_descent(f, grad, x0, step_rule, tol=1e-8, max_iter=25_000):
    x = np.asarray(x0, dtype=float).copy()
    xs, values, norms, steps = [], [], [], []
    for _ in range(max_iter + 1):
        value, g = float(f(x)), np.asarray(grad(x), dtype=float)
        xs.append(x.copy()); values.append(value)
        norms.append(np.linalg.norm(g))
        if norms[-1] <= tol:
            break
        alpha = float(step_rule(f, x, g))
        steps.append(alpha)
        x = x - alpha * g
    return {"x": x, "iterates": np.asarray(xs),
            "values": np.asarray(values),
            "gradient_norms": np.asarray(norms),
            "step_sizes": np.asarray(steps)}

def armijo_step(f, x, g, initial=1.0, c=1e-4, rho=0.5):
    alpha, fx = initial, f(x)
    while f(x - alpha*g) > fx - c*alpha*(g @ g):
        alpha *= rho
    return alpha
\end{lstlisting}

For the quadratic experiment, seed \(3\) gives
\[
\lambda_{\min}=0.115303501,
\quad \lambda_{\max}=74.677637749,
\quad \kappa(Q)=647.661.
\]
The benchmark is computed with
\texttt{np.linalg.solve(Q, c)}, never by explicitly forming \(Q^{-1}\).

\begin{center}
\begin{tabular}{lrrrl}
\toprule
Step & Iterations & Final \(\norm{\nabla f}\) & \(\norm{x-x^\star}\) & Status\\
\midrule
\(1/\lambda_{\max}\) & 11199 & \(9.99\times10^{-9}\) & \(8.66\times10^{-8}\) & converged\\
\(2/(\lambda_{\min}+\lambda_{\max})\) & 5656 & \(9.99\times10^{-9}\) & \(7.37\times10^{-8}\) & converged\\
\(2.05/\lambda_{\max}\) & 7358 & overflow & --- & unstable\\
\bottomrule
\end{tabular}
\end{center}
The second stable step roughly halves the iteration count. The third violates
the stability bound and eventually amplifies the largest-eigenvalue component.

For Rosenbrock,
\[
\nabla r(x,y)=
\binom{-400x(y-x^2)-2(1-x)}{200(y-x^2)}.
\]
At \((-1.2,1)\), a centered finite-difference check with \(h=10^{-6}\) has
infinity-norm error \(2.24\times10^{-8}\). Armijo backtracking with
\(c=10^{-4}\), \(\rho=1/2\), and initial trial step \(1\) reaches
\[
x=(0.99999999,0.99999998),
\quad r(x)=6.11\times10^{-17},
\quad \norm{\nabla r(x)}=1.00\times10^{-8}
\]
after 19435 accepted steps. Slow motion along the narrow curved valley is
expected for plain steepest descent.

The reference script produces the four diagnostics requested in the exercise:
the quadratic objective gap, the quadratic gradient norm, the path in the
extreme eigendirections of \(Q\), and the Rosenbrock path over contours. A
correct plot uses a positive floor before taking a logarithm of an objective
gap at machine precision.

\paragraph{Exercise 11 -- Implement primal--dual gradient descent--ascent.}

The exact reference pair is obtained from
\[
\begin{pmatrix}Q&A\trans\\A&0\end{pmatrix}
\binom{x^\star}{\lambda^\star}=\binom{-c}{b}.
\]
For \(e_k=x_k-x^\star\) and \(d_k=\lambda_k-\lambda^\star\), the simultaneous
iteration has transition matrix
\[
\binom{e_{k+1}}{d_{k+1}}
=
\underbrace{\begin{pmatrix}
I-\alpha Q&-\alpha A\trans\\
\beta A&I
\end{pmatrix}}_{T_{\alpha,\beta}}
\binom{e_k}{d_k}.
\]
Thus the linear iteration converges for every initial condition exactly when
\(\rho(T_{\alpha,\beta})<1\). This spectral-radius calculation explains the
observed behavior more precisely than a visual judgment from the trajectories.

Using seed \(11\), \(d=20\), \(m=3\), and the construction in the exercise:
\[
\begin{array}{c|c|c|r|c}
(\alpha,\beta)&\rho(T)&\text{behavior}&\text{iterations}&
\max(\text{final residuals})\\ \midrule
(0.02,0.02)&0.988187&\text{stable}&1621&9.92\times10^{-9}\\
(0.03,0.03)&1.064405&\text{growing oscillations}&3038&\text{diverged}
\end{array}
\]
The Lagrangian value need not decrease monotonically: one variable descends
while the multiplier ascends. Feasibility and stationarity residuals, not the
Lagrangian alone, certify convergence.

For the augmented Lagrangian,
\[
\nabla_x\mathcal L_\rho
=Qx+c+A\trans\lambda+\rho A\trans(Ax-b),
\]
so the primal block in the transition matrix becomes
\(I-\alpha(Q+\rho A\trans A)\). With
\((\alpha,\beta)=(0.005,0.02)\), the reproducible comparison is
\[
\begin{array}{c|c|r|cc}
\rho&\text{spectral radius}&\text{iterations}&
\text{feasibility}&\text{stationarity}\\ \midrule
0.1&0.996814&5770&6.84\times10^{-10}&9.98\times10^{-9}\\
1&0.996803&5766&7.70\times10^{-10}&9.99\times10^{-9}\\
10&0.998026&9916&2.40\times10^{-10}&1.00\times10^{-8}
\end{array}
\]
Larger \(\rho\) penalizes infeasibility more strongly but also increases primal
curvature, so a step size that is effective at \(\rho=1\) need not remain
effective at \(\rho=10\). All converged solutions agree with the exact KKT
solution to the reported tolerance.

For the optional extragradient method, define
\[
F(x,\lambda)=
\binom{Qx+c+A\trans\lambda}{-(Ax-b)}.
\]
Then compute
\[
\widetilde z_k=z_k-\eta F(z_k),
\qquad
z_{k+1}=z_k-\eta F(\widetilde z_k).
\]
The look-ahead evaluation damps the rotational behavior created by the
off-diagonal \(A,A\trans\) blocks and often permits a wider stable step range
than simultaneous descent--ascent. Stability should still be verified from
residual histories or the corresponding linear iteration matrix.

% -----------------------------------------------------------------------------
\section{Portfolio Optimization}

\paragraph{Exercise 12 -- Risk, return, and the Sharpe ratio.}

For deterministic weights \(w\), linearity of expectation and covariance give
\[
\mathbb E[w\trans R]=w\trans\mu,
\qquad
\operatorname{Var}(w\trans R)
=\mathbb E[(w\trans(R-\mu))^2]=w\trans\Sigma w.
\]

For GMV, use
\(\mathcal L(w,\lambda)=\frac12w\trans\Sigma w
+\lambda(\mathbf1\trans w-1)\). Stationarity gives
\(w=-\lambda\Sigma^{-1}\mathbf1\), and the budget fixes the multiplier.
Thus
\begin{resultbox}
\[
w_{\mathrm{GMV}}
=\frac{\Sigma^{-1}\mathbf1}
{\mathbf1\trans\Sigma^{-1}\mathbf1}.
\]
\end{resultbox}

Let \(e=\mu-r_f\mathbf1\). Cauchy--Schwarz in the \(\Sigma\)-inner product
gives
\[
e\trans w
=(\Sigma^{-1/2}e)\trans(\Sigma^{1/2}w)
\leq\sqrt{e\trans\Sigma^{-1}e}\sqrt{w\trans\Sigma w}.
\]
Equality holds precisely in the tangency direction
\(w\propto\Sigma^{-1}e\). If
\(D=\mathbf1\trans\Sigma^{-1}e\ne0\), full investment fixes
\[
w_T=\frac{\Sigma^{-1}(\mu-r_f\mathbf1)}{D}.
\]
For the usual maximum-positive-Sharpe interpretation one also checks the sign
of \(D\); a negative normalization reverses the direction.

Equal weight is a transparent, low-estimation benchmark. GMV uses covariance
but not expected returns and is appropriate when risk control is primary.
Tangency targets risk-adjusted excess return but is especially sensitive to
mean estimates. With long-only constraints the closed forms can contain
negative weights and become infeasible. For long-only GMV, for example, KKT
requires
\[
\Sigma w+\lambda\mathbf1-\nu=0,
\quad \mathbf1\trans w=1,
\quad w\geq0,
\quad \nu\geq0,
\quad \nu_iw_i=0.
\]
An analogous stationarity equation, together with the same primal, dual, and
complementarity conditions, is required for the chosen tangency formulation.

\paragraph{Exercise 13 -- Derive the efficient frontier.}

Let
\[
C=\begin{pmatrix}\mathbf1\trans\\\mu\trans\end{pmatrix},
\qquad d_m=\binom1m.
\]
The Lagrangian gives the KKT system
\[
\begin{pmatrix}\Sigma&C\trans\\C&0\end{pmatrix}
\binom{w}{\lambda}=\binom0{d_m}.
\]
Eliminating \(w\), and assuming the two constraints are independent, yields
\begin{resultbox}
\[
w(m)=\Sigma^{-1}C\trans
(C\Sigma^{-1}C\trans)^{-1}d_m.
\]
\end{resultbox}
Define
\[
A_0=\mathbf1\trans\Sigma^{-1}\mathbf1,
\quad B_0=\mathbf1\trans\Sigma^{-1}\mu,
\quad D_0=\mu\trans\Sigma^{-1}\mu,
\quad \Delta=A_0D_0-B_0^2>0.
\]
Since the multiplier calculation gives
\(w(m)\trans\Sigma w(m)=d_m\trans
(C\Sigma^{-1}C\trans)^{-1}d_m\), the frontier variance is
\begin{resultbox}
\[
\sigma^2(m)=\frac{A_0m^2-2B_0m+D_0}{\Delta}.
\]
\end{resultbox}
It is a convex quadratic in target return. In mean--standard-deviation space
it produces a hyperbola; the efficient branch begins at the GMV return
\(m_{\mathrm{GMV}}=B_0/A_0\) and moves toward higher returns.

Long-only constraints and caps make the feasible set polyhedral. As \(m\)
changes, weights hit or leave bounds; each fixed active set has its own
quadratic formula, producing piecewise behavior and visible active-set
changes. Before tracing a frontier, code should verify symmetry and positive
semidefiniteness of \(\Sigma\), preferably positive definiteness for
uniqueness; rank and consistency of the equality constraints; feasibility of
the budget and target; bounds and caps; solver success; and final KKT
residuals. Under long-only full investment, a necessary and sufficient return
check without caps is \(m\in[\min_i\mu_i,\max_i\mu_i]\).

\paragraph{Exercise 14 -- Estimation error and robust weights.}

Write
\(q=\widehat\Sigma^{-1}(\widehat\mu-r_f\mathbf1)\). Keeping covariance and
the risk-free rate fixed gives the exact perturbation formula
\[
q(\widehat\mu+\varepsilon h,\widehat\Sigma)
=q(\widehat\mu,\widehat\Sigma)
+\varepsilon\widehat\Sigma^{-1}h.
\]
Therefore
\[
D_\mu q[h]=\widehat\Sigma^{-1}h,
\qquad J_\mu q=\widehat\Sigma^{-1}.
\]
If
\(\widehat\Sigma=U\diag(\lambda_1,\ldots,\lambda_d)U\trans\) and
\(h=\sum_i h_i u_i\), then
\[
\Delta q=\varepsilon\sum_i\frac{h_i}{\lambda_i}u_i,
\qquad
\norm{\Delta q}_2
\leq\frac{|\varepsilon|}{\lambda_{\min}(\widehat\Sigma)}\norm h_2.
\]
Small-eigenvalue directions amplify mean error the most. The corresponding
relative bound is the right-hand side divided by
\(\norm q_2\), when \(q\ne0\).

For normalized weights \(w=q/(\mathbf1\trans q)\), set
\(D=\mathbf1\trans q\). Differentiation gives
\[
D_\mu w[h]
=\frac{D\widehat\Sigma^{-1}h
-q\,\mathbf1\trans\widehat\Sigma^{-1}h}{D^2}.
\]
The denominator explains the additional instability when \(D\) is close to
zero.

GMV weights do not contain \(\widehat\mu\), so they are immune to mean
estimation error. They still contain both \(\widehat\Sigma^{-1}\) and a
normalizing denominator, and hence remain sensitive to covariance error,
especially near singularity.

For ridge regularization, the eigenvectors are unchanged and
\[
\lambda_i(\widehat\Sigma_\delta)=\lambda_i+\delta,
\qquad
\kappa(\widehat\Sigma_\delta)
=\frac{\lambda_{\max}+\delta}{\lambda_{\min}+\delta}.
\]
This condition number decreases toward \(1\) as \(\delta\) increases, although
the estimator becomes more biased toward an isotropic covariance.

In-sample predicted volatility reuses the observations that selected both the
covariance and the weights, so it is optimistically biased. Out-of-sample
realized volatility must be computed from untouched future observations.
Expected returns are particularly noisy, and optimization applies an inverse
covariance to that noise; the resulting extreme weights can make return-based
objectives much less stable than GMV.

A transparent selection rule is:
\begin{enumerate}
  \item split the training period chronologically into an inner estimation
  sample and a validation sample;
  \item declare a grid such as
  \(\delta=\tau\operatorname{tr}(\widehat\Sigma)/d\) before inspecting the
  final test set;
  \item fit each candidate on the inner sample and rank it by validation
  realized variance, optionally with a predeclared concentration or turnover
  tie-breaker;
  \item choose \(\delta\), refit using all training observations, and evaluate
  the final test sample exactly once.
\end{enumerate}
Rolling validation may replace a single split, provided every validation fold
occurs strictly after its estimation window and the final test set remains
untouched.

% -----------------------------------------------------------------------------
\begin{projectbox}{Separate Portfolio Optimization Project}
The correction to the separate Portfolio Optimization Project is still being
drafted. It is intentionally not included in this correction document.
\end{projectbox}

\end{document}
